% Options for packages loaded elsewhere
\PassOptionsToPackage{unicode}{hyperref}
\PassOptionsToPackage{hyphens}{url}
\documentclass[
  11pt,
  a4paper,
]{article}
\usepackage{xcolor}
\usepackage[margin=25mm]{geometry}
\usepackage{amsmath,amssymb}
\setcounter{secnumdepth}{-\maxdimen} % remove section numbering
\usepackage{iftex}
\ifPDFTeX
  \usepackage[T1]{fontenc}
  \usepackage[utf8]{inputenc}
  \usepackage{textcomp} % provide euro and other symbols
\else % if luatex or xetex
  \usepackage{unicode-math} % this also loads fontspec
  \defaultfontfeatures{Scale=MatchLowercase}
  \defaultfontfeatures[\rmfamily]{Ligatures=TeX,Scale=1}
\fi
\usepackage{lmodern}
\ifPDFTeX\else
  % xetex/luatex font selection
\fi
% Use upquote if available, for straight quotes in verbatim environments
\IfFileExists{upquote.sty}{\usepackage{upquote}}{}
\IfFileExists{microtype.sty}{% use microtype if available
  \usepackage[]{microtype}
  \UseMicrotypeSet[protrusion]{basicmath} % disable protrusion for tt fonts
}{}
\makeatletter
\@ifundefined{KOMAClassName}{% if non-KOMA class
  \IfFileExists{parskip.sty}{%
    \usepackage{parskip}
  }{% else
    \setlength{\parindent}{0pt}
    \setlength{\parskip}{6pt plus 2pt minus 1pt}}
}{% if KOMA class
  \KOMAoptions{parskip=half}}
\makeatother
\usepackage{longtable,booktabs,array}
\usepackage{caption}
\captionsetup[table]{skip=6pt}
\newcounter{none} % for unnumbered tables
\usepackage{calc} % for calculating minipage widths
% Correct order of tables after \paragraph or \subparagraph
\usepackage{etoolbox}
\makeatletter
\patchcmd\longtable{\par}{\if@noskipsec\mbox{}\fi\par}{}{}
\makeatother
% Allow footnotes in longtable head/foot
\IfFileExists{footnotehyper.sty}{\usepackage{footnotehyper}}{\usepackage{footnote}}
\makesavenoteenv{longtable}
\usepackage{graphicx}
\makeatletter
\newsavebox\pandoc@box
\newcommand*\pandocbounded[1]{% scales image to fit in text height/width
  \sbox\pandoc@box{#1}%
  \Gscale@div\@tempa{\textheight}{\dimexpr\ht\pandoc@box+\dp\pandoc@box\relax}%
  \Gscale@div\@tempb{\linewidth}{\wd\pandoc@box}%
  \ifdim\@tempb\p@<\@tempa\p@\let\@tempa\@tempb\fi% select the smaller of both
  \ifdim\@tempa\p@<\p@\scalebox{\@tempa}{\usebox\pandoc@box}%
  \else\usebox{\pandoc@box}%
  \fi%
}
% Set default figure placement to htbp
\def\fps@figure{htbp}
\makeatother
\setlength{\emergencystretch}{3em} % prevent overfull lines
\providecommand{\tightlist}{%
  \setlength{\itemsep}{0pt}\setlength{\parskip}{0pt}}
\usepackage{amsmath,amssymb}
\usepackage{xurl}
\usepackage{fvextra}
\DefineVerbatimEnvironment{verbatim}{Verbatim}{breaklines=true,fontsize=\small}
\urlstyle{same}
\setlength{\emergencystretch}{2em}
\AtBeginDocument{\hypersetup{colorlinks=true,linkcolor=blue,urlcolor=blue}}

% Compatibility: Pandoc image alt metadata with the retained TeX Live 2020.
\makeatletter
\define@key{Gin}{alt}{}
\makeatother
\usepackage{bookmark}
\IfFileExists{xurl.sty}{\usepackage{xurl}}{} % add URL line breaks if available
\urlstyle{same}
% fallback for those not using the hyperref driver hyperxmp:
\makeatletter
\@ifundefined{xmpquote}{\newcommand{\xmpquote}[1]{#1}}{}
\makeatother
\hypersetup{
  pdftitle={Minimum Widths for Unit Squares in Thin Rectangles},
  pdfauthor={Hiroaki Hosono (working draft)},
  hidelinks,
  pdfcreator={LaTeX via pandoc}}

\title{Minimum Widths for Unit Squares in Thin Rectangles}
\author{Hiroaki Hosono (working draft)}
\date{25 October 2026}

\begin{document}
\maketitle

\section{Abstract}\label{abstract}

We present a computer-assisted determination of the minimum width needed
to pack \(N\) unit squares, with arbitrary rotations and pairwise
disjoint interiors, in a rectangle of height \(1851/1000\le a<2\). For
odd \(N\ge11\), the width is \(4+2(N-2)\sqrt2-(N-1)a\). For even
\(N\ge12\), it is the minimum of an explicit rationally parametrised
function, attained by a staggered configuration with two upright end
caps. The representative minimising angle is unique. The proof separates
a finite all-sequence type-forcing certificate from analytical
minimisation and attainment. We describe the exact acceptance
obligations and the retained two-system evidence. The result does not
determine the entire rectangle frontier for 16--20 squares.

\section{1. Introduction and notation}\label{introduction-and-notation}

Let \(\nu(x,a)\) be the maximum number of closed unit squares contained
in \([0,x]\times[0,a]\) with pairwise disjoint interiors, allowing
arbitrary orientations. Contacts are permitted. Write

\[
B_N(a)=\inf\{x>0:\nu(x,a)\ge N\}.
\]

Friedman's rectangle problem {[}1{]} records limiting packings and
boundary pieces for small numbers of squares. Here it supplies
construction history, not a certificate excluding every smaller packing.
Nagamochi {[}2{]} studies broader rectangle counting bounds with nominal
side parameters at least two. Karakuş's preprint {[}3{]} reports a
defect in the scoring assertion of that proof and provides a separate
weaker rectangle bound. Neither scoring argument is used here.

Our question fixes a height strictly below two and determines the width
uniformly in the number of squares. The proof below does not depend on
the separate strip and strict zipper-corner lemmas studied elsewhere in
this project.

The search for better 16--20-square constructions at other heights is a
separate upper-bound project. Its tables are not proof of the complete
feasible frontier and are not used to establish the theorem below.

No priority claim is made. The references were checked at the levels
described in Section 8; this is not an exhaustive literature search.

\section{2. Main theorem}\label{main-theorem}

Represent an orientation by \(t=\tan(\theta/2)\) and let

\[
\mathcal O=\left[2\arctan\frac{4339}{12800},\;2\arctan\frac{4553}{8800}\right].
\]

For \(\theta\in\mathcal O\), put \(c=\cos\theta\), \(s=\sin\theta\),
\(h=c+s\), \(b=3a/2-1\), and

\[
P=\frac{2h-a}{h^2-1},\qquad
Q=1+\frac{2h-b}{h^2-1},\qquad
A=\frac{2-a}{4}(\tan\theta-\cot\theta).
\]

Define

\[
G_N(\theta;a)=(N-4)P+2Q-2|A|.
\]

All denominators are positive on this angle interval. The symbol \(h\)
denotes a square's bounding-box side length, whereas \(a\) is the
container height.

\textbf{Theorem 2.1 (exact thin-strip width; computer-assisted).} For
every \(1851/1000\le a<2\):

\begin{enumerate}
\def\labelenumi{\arabic{enumi}.}
\tightlist
\item
  For odd \(N\ge11\), \[
  B_N(a)=4+2(N-2)\sqrt2-(N-1)a.
  \]
\item
  For even \(N\ge12\), \[
  B_N(a)=\min_{\theta\in\mathcal O}G_N(\theta;a).
  \] The representative minimiser \(\theta_N(a)>\pi/4\) is unique and
  lies below \(2\arctan(43/100)\). Its reflection \(\pi/2-\theta_N(a)\)
  gives the other minimising orientation. The value is attained by a
  staggered packing.
\end{enumerate}

The theorem combines the retained two-system results on
\([1851/1000,19/10]\) and \([19/10,2)\), including the separate coverage
for \(N=11\). The mathematical argument is given below, with its finite
proof obligations exposed in Section 6. The new manuscript is not itself
a fresh replay of those computations.

\section{3. Certified type forcing}\label{certified-type-forcing}

Put \(g=2\sqrt2-a\) and \(F_0=4-2\sqrt2\). The retained certificate
establishes the following assertion for every \(N\ge11\): a packing of
width at most

\[
(N-1)g+F_0
\]

has two near-upright end caps, its \(N-2\) interior squares have angles
in \(\mathcal O\), and the centres of the interior squares alternate
between the closed lower and upper halves of the container (the midpoint
belongs to both). Every non-type sequence has strictly greater width.

This is an all-sequence assertion. The verification covers angle/height
cells and all transitions, not just configurations found in a search.
Exact separating-axis bounds and rational potentials control non-type
transitions. Calibrated type transitions use angular variation. A
saturated dynamic-programming layer and its checked fixed point cover
arbitrarily large \(N\).

The threshold is sufficient for the equality theorem as follows. The
\(45^\circ\) construction has width \((N-1)g+F_0\). For even \(N\),
\(G_N(\pi/4;a)\) equals that threshold, so \(\min G_N\) is no larger. An
arbitrary packing wider than the threshold already satisfies the claimed
lower bound. Any packing at or below the threshold is forced into the
type, where Section 4 applies. Section 5 constructs a packing attaining
the resulting value. Thus type forcing alone is not being substituted
for minimisation or attainment.

The finite verification includes transitions into a type sequence, not
only transitions within one. Near height two it uses a corrected
absolute-value inequality for nearly upright squares. Earlier
computations with the superseded signed inequality are excluded from the
evidence.

\section{4. The lower bound inside the
type}\label{the-lower-bound-inside-the-type}

\subsection{4.1. Pair and cap
inequalities}\label{pair-and-cap-inequalities}

For interior angles \(\theta_1,\ldots,\theta_{N-2}\) define

\[
\mathrm{TV}=\sum_{i=1}^{N-3}|\theta_{i+1}-\theta_i|.
\]

The retained pair calibration gives a lower bound per radian of
\(\theta\). Its constant can be chosen as

\[
k=\frac15\quad\left(\frac{1851}{1000}\le a\le\frac{19}{10}\right),
\qquad
k=\frac{227}{2000}\quad\left(\frac{19}{10}\le a<2\right).
\]

Together with the cap comparison, summing the inequalities produces one
\(P\)-term for each non-end interior square, end-interior terms
\(Q\pm A\), and the nonnegative contribution \(k\,\mathrm{TV}\). The
same certified bounds give \(|\partial_\theta A|\le k\).

The cap comparison is applied to pairs at their original positions. It
does not assume that simultaneously moving arbitrary caps preserves a
packing. The lower-height cap comparison is a separate certified
dependency; the upper-height argument must not be extrapolated below 1.9
without it.

\subsection{4.2. Reduction to a common
angle}\label{reduction-to-a-common-angle}

Choose an interior angle \(\bar\theta\) closest to \(\pi/4\). Both \(P\)
and \(Q\) decrease as \(h\) increases. Indeed,

\[
\frac{d}{dh}\frac{2h-z}{h^2-1}
=-\frac{2(h^2-zh+1)}{(h^2-1)^2}<0
\]

for \(z=a\) or \(z=b\), since \(z<2\) in the claimed range. Consequently
each \(P\)- and \(Q\)-term is at least its value at \(\bar\theta\).

For even \(N\), the two endpoint \(A\)-terms have the same sign pattern.
Replacing them by \(-2|A(\bar\theta)|\) loses at most

\[
k\bigl(|\theta_1-\bar\theta|+|\theta_{N-2}-\bar\theta|\bigr)
\le k\,\mathrm{TV}.
\]

The variation contribution absorbs this loss. The occupied width is
therefore at least \(G_N(\bar\theta;a)\), and hence at least its
minimum.

For odd \(N\), the endpoint \(A\)-terms have opposite signs. Their
difference is absorbed by \(k\,\mathrm{TV}\). The remainder is minimised
when \(h=\sqrt2\), giving

\[
(N-4)P(\pi/4)+2Q(\pi/4)
=4+2(N-2)\sqrt2-(N-1)a.
\]

\subsection{4.3. Uniqueness and the critical
equation}\label{uniqueness-and-the-critical-equation}

In this paragraph differentiation is with respect to \textbf{\(t\)}, not
per radian of \(\theta\). On

\[
I=[\sqrt2-1,4553/8800],
\]

the retained exact polynomial checks establish

\[
P_{tt}>0,\qquad P_{tt}+2Q_{tt}-2A_{tt}>0.
\]

On the upper-angle side \(A\ge0\), and

\[
G_{N,tt}=(N-5)P_{tt}+(P_{tt}+2Q_{tt}-2A_{tt})>0.
\]

The derivative at \(t=\sqrt2-1\) is

\[
G_{N,t}=-(2-a)(2+\sqrt2)<0,
\]

while the checked derivative at \(t=43/100\) is positive for \(N\ge12\).
This yields a unique representative critical point below that endpoint.
The reflection of the lower-angle part of \(\mathcal O\) lies in the
checked upper interval, so it adds no smaller value. Reflection gives
the other minimising orientation.

The rational parametrisation of \(c,s\) turns the critical-point
equation into a polynomial equation in \(t\), with rational polynomial
coefficients in \(a\), after multiplication by positive denominators.
The minimising parameter is therefore an algebraic function of \(a\).
When \(a\) is algebraic, its parameter and minimum width are algebraic
numbers. For an arbitrary real \(a\), no algebraicity over the rationals
is asserted, and the angle \(2\arctan t\) is not being called an
algebraic number. Convexity and endpoint signs here are certificate
dependencies, not deductions from a plot.

\section{5. Attaining configurations}\label{attaining-configurations}

Use two upright caps and \(N-2\) interiors at one angle, on alternating
walls. Let

\[
f_+=\frac{1-s(a-h)}c,\qquad
f_-=\frac{1-c(a-h)}s,
\]

\[
c_+=\frac{(1+h)/2-s(a-(1+h)/2)}c,\qquad
c_-=\frac{(1+h)/2-c(a-(1+h)/2)}s.
\]

Interior centres have height \(h/2\) or \(a-h/2\). Each cap has centre
height \(1/2\) or \(a-1/2\), on the wall opposite its neighbour.
Starting with cap centre \(x=1/2\), take consecutive horizontal
increments \(c_\pm\) at the caps and \(f_\pm\) between interiors,
according to the alternating wall pattern.

Adjacent pairs touch on an appropriate face normal. On the angle
interval containing every minimiser, the retained attainment checks
establish

\[
f_\pm>h/2,\qquad
c_\sigma+f_{\sigma'}>(1+h)/2,\qquad
c_\sigma>(h-1)/2
\]

for the required sign choices. The first two inequalities separate
nonadjacent pairs horizontally; longer chains add positive increments.
The last makes the caps outermost. The cap--cap pair is also separated,
since a cap increment and two interior increments already exceed one.

Containment follows from \(h\le\sqrt2<a\). Thus the width is exactly one
plus the sum of the horizontal increments. For even \(N\), choose the
starting wall to obtain

\[
1+2c_{\sigma_0}+f_{-\sigma_0}
 +(N-4)(f_++f_-)/2
=(N-4)P+2Q-2|A|.
\]

For odd \(N\), use \(\theta=\pi/4\); the same sum is the formula in
Theorem 2.1. Choosing the critical orientation in the even case gives
exact attainment. The alternating sequence puts the two upright caps on
the same wall for odd numbers of squares and on opposite walls for even
numbers (Figure 1).

For fixed \(t\), the individual expressions in the analysis and
attainment checks are affine in \(a\). The accepted endpoint
certificates therefore extend pointwise through the two height
intervals. This is applied to individual inequalities, not to a presumed
endpoint extremum of an arbitrary minimum of affine branches.

A nearby rational-half-angle placement certificate may have a small
positive width slack. Such a certificate is evidence of a packing near
the optimum, not equality of its width with the critical value.

\par\medskip
\noindent\begin{minipage}{\linewidth}
\centering
\pandocbounded{\includegraphics[keepaspectratio,alt={Alternating-wall constructions for odd and even numbers of squares.}]{figures/parity-packing.pdf}}

\par\smallskip\raggedright
\emph{Figure 1. The construction of Section 5 at height 1.9: eleven
squares at 45 degrees (the odd optimum), and twelve squares at the
illustrative rational half-angle parameter 427/1000. Orange squares are
the two upright end caps; blue squares alternate between the walls. The
even panel illustrates a feasible construction at the displayed angle,
not the exact critical-angle optimum. Coordinates are rounded only for
drawing; the figure is not a packing certificate.}
\end{minipage}\par\medskip

\par\medskip
\noindent\begin{minipage}{\linewidth}
\centering
\pandocbounded{\includegraphics[keepaspectratio,alt={Height-width formulas for eleven squares and illustrative twelve-square constructions.}]{figures/height-width.pdf}}

\par\smallskip\raggedright
\emph{Figure 2. Height and width on the theorem interval. The odd curve
is the exact formula for eleven squares. The even curve is the width of
the twelve-square construction at the fixed half-angle parameter
427/1000; it is an upper bound, while the theorem takes the minimum over
all admissible angles. The right endpoint at height two is excluded.
Lines are drawn from the stated formulas and do not establish the finite
lower-bound computation.}
\end{minipage}\par\medskip

\section{6. Exact computations and
reproducibility}\label{exact-computations-and-reproducibility}

\subsection{6.1. Retained evidence}\label{retained-evidence}

All identifiers below are relative to rect\_search/claude/g2. These are
existing internal receipts, not new checks carried out for this paper or
an external peer review.

{\def\LTcaptype{none} % do not increment counter
\begin{longtable}[]{@{}
  >{\raggedright\arraybackslash}p{(\linewidth - 2\tabcolsep) * \real{0.5000}}
  >{\raggedright\arraybackslash}p{(\linewidth - 2\tabcolsep) * \real{0.5000}}@{}}
\toprule\noalign{}
\begin{minipage}[b]{\linewidth}\raggedright
Ingredient
\end{minipage} & \begin{minipage}[b]{\linewidth}\raggedright
Specification or receipt
\end{minipage} \\
\midrule\noalign{}
\endhead
\bottomrule\noalign{}
\endlastfoot
state cover, exact bounds, DP and saturation rules &
STAGE1\_DP\_SPEC.md \\
upper-height analysis and attainment & STAGE2\_INTERVAL.md \\
upper-height evidence, including the corrected final interval and
\(N=11\) coverage & EVIDENCE\_INDEX.md \\
lower-height type forcing, repaired transitions and saturation &
s4e/RECEIPT.md \\
lower-height calibration and cap comparison & s4b/CAL\_LOW\_SPEC.md;
s4b/CAP\_UPRIGHT\_LOW\_SPEC.md \\
lower-height analysis and attainment & s4b/STAGE2\_LOW.md;
s4b/STAGE2\_LOW\_RECEIPT.md \\
upper-height calibration & CAL\_INTERVAL\_SPEC.md \\
\end{longtable}
}

The final lower-height verification record, dated 21 October 2026,
documents checks by the first implementation and a second implementation
derived from the specifications, followed by internal mathematical
review. It reports complete interval and Sturm checks separately from
the 36 finite placement tests. Those placements are not a proof of all
heights or all \(N\).

For \(a\in[1851/1000,1852/1000]\), the s4g retake establishes \(N\ge24\)
using its individual finite layers together with its saturated tail; the
first-system output labels that tail \(N\ge26\), rather than \(N\ge24\).
Separate s3q fixed-\(N\) runs cover \(11\le N\le26\). Their union covers
every \(N\ge11\). On \([1852/1000,1854/1000]\), the chord runs cover
\(N\ge12\) and separate runs cover \(N=11\). Later-height runs cover the
remainder. The corrected absolute-form final interval and the separate
\(N=11\) subdivision at height 1.9 must be included.

The second implementation was written from specifications rather than
the first implementation's code, as recorded by its read log and
receipts. That independence is a retained provenance claim, not a new
independent audit performed during manuscript preparation.

\subsection{6.2. Acceptance obligations}\label{acceptance-obligations}

A complete replay must check the closed state cover, all SAT branches,
potential coverage, every transition (including entry to and exit from a
type run), the integer rounding model, and the saturated-layer fixed
point. It must also check the full-domain calibration and cap
inequalities, the polynomial signs and root counts, the affinity
argument and the attainment inequalities. Positive output on selected
intervals, floating LP feasibility, placement samples or a root count
alone are insufficient.

\subsection{6.3. Retained command
interfaces}\label{retained-command-interfaces}

The following examples use the working directory rect\_search/claude/g2.
They are \textbf{instructions, not commands executed for this
manuscript}.

\begin{verbatim}
python3 s4b/stage2_low.py 1851/1000 19/10
python3 s4b/stage2_low_affine.py
python3 check2_9b/stage2_iv_check2.py --a-lo 1851/1000 --a-end 19/10
python3 s4b/attain_low.py 43/100
python3 check2_9b/attain_check2.py --a-lo 1851/1000 --a-hi 19/10
python3 check2_9b/cal_low_check2.py 1851/1000
python3 check2_9b/cal_low_check2.py 19/10
python3 check2_9b/cap_upright_check2.py 1024 128 --hcap1
python3 check2_9b/stage2_iv_check2.py
python3 check2_9b/attain_check2.py
\end{verbatim}

These source interfaces cover analysis, calibration and attainment; they
are not an exhaustive type-forcing replay. For that part the retained
chord checker takes interval endpoints, state/certificate files and an
output path, with proof-variant switches:

\begin{verbatim}
python3 check2_9b/g2s1_chord.py ALO AHI STATES CERT --out OUT
    [--all] [--phi0] [--nosat] [--L N] [--uu]
\end{verbatim}

The accepted interval/input lists and result headers are indexed in the
evidence documents. In particular, fixed-\(N\) no-saturation runs do not
replace the separately required saturation check. The final interval
\([1.999,2)\) requires the second checker's near-upright switch,
\texttt{-\/-uu}, with the accepted absolute-form rule; the switch is
read when its interval object is constructed.

Python 3, SymPy and NumPy occur in the retained implementations. Before
submission, a frozen portable supplement must pin their actual
dependency/environment records, corrected production code and all
accepted input/output identities. Missing versions are not invented
here. No recorded hash was recomputed and no new public supplement was
assembled in this task.

\section{7. Limits and remaining work}\label{limits-and-remaining-work}

The endpoint \(a=1851/1000\) is included; \(a=2\) is excluded. The
theorem does not include \(N\le10\). The strict zipper-corner theorem
from the other workstream does not fill the missing height-two boundary.

The retained lower-end gap is approximately

\[
a_{11}\le a<1851/1000,\qquad
a_{11}=\frac{18\sqrt2-7}{10}\approx1.8455844.
\]

This expression marks the intersection of the odd eleven-square
staggered line with width eleven; it is not a proved common transition
for every \(N\). Extending the all-packing theorem below 1.851 requires
additional accepted proof coverage. A negative DP lower bound is
inconclusive, not a counterexample packing.

The full frontier for \(N=16,17,18,19,20\) at other heights remains a
separate problem. Construction tables and packing certificates establish
upper bounds only; they do not exclude smaller packings.

This version is a working paper, with TeX and PDF companions, and has
not been submitted to a journal or externally peer reviewed. A complete
portable computational supplement, including its environment and replay
inventory, remains to be assembled. The assembled mathematical text has
received a limited independent consistency review against the retained
records; that review is not a fresh execution of the proof computations.
A broader literature and priority search, and final authorship and
acknowledgements, also remain. No mathematical checker, solver,
placement test or certificate replay was run in preparing this text or
its figures.

\section{8. References and verification
level}\label{references-and-verification-level}

{[}1{]} Erich Friedman, \emph{Problem of the Month (June 2010)}, Math
Magic. \url{https://erich-friedman.github.io/mathmagic/0610.html}.
Primary problem page and displayed construction table inspected on 25
October 2026.

{[}2{]} Hiroshi Nagamochi, \emph{Packing Unit Squares in a Rectangle},
The Electronic Journal of Combinatorics \textbf{12} (2005), Research
Paper R37. \url{https://doi.org/10.37236/1934}. Journal bibliographic
record and abstract inspected; its proof is not reused.

{[}3{]} Hakan Karakuş, \emph{A counterexample to Nagamochi's scoring
lemma and a new rectangle packing bound}, arXiv:2609.37410v1 (2026).
\url{https://arxiv.org/abs/2609.37410v1}. Author, title, submission date
and abstract inspected; the full proof was not audited here.

Internal source identifiers for the proposed computational supplement
are listed in Section 6. They are project evidence records, not
fabricated public bibliographic entries.

\end{document}
